\section{Limitations and Future Work}
\label{sec:limitations}

\subsection{Evidence gaps}
The largest limitation is the absence of a controlled study. This paper validates the implementation, the loop's transitions, and a nine-document corpus, but it does not show that a lesson improves a paper, a grant, or an analysis. The corpus is offline and drawn from one user's machine; its audits are cold-start overlays, so the finding mix reflects missing setup as well as manuscript problems. The next step is a pre-registered live evaluation with three properties. First, reserve entire projects or task batches for reporting and never use them to propose or admit lessons. Second, run both arms with the same model, tools, prompt hashes, and inference budget, and repeat each task enough times to estimate variation. Third, use blinded human review for quality and a deterministic audit for hard constraints. The survey's Selection, Activation, Attribution, and Improvement rubric is the right reporting frame \citep{wei2026harnessrsi}; the current system records the episodes and decisions that such a study would need.

\subsection{Outcome and activation}
The outcome score is a heuristic. It weights errors, corrections, blocking findings, integrity, citation, statistics, figures, reproducibility, gate state, and compile status, but the weights are not calibrated. A useful extension is to separate the score into a hard-constraint vector and a quality estimate, and to report both rather than a single number. The loop should also measure task completion and coverage: brief fulfillment, inventory coverage, reviewer-item resolution, and unnecessary handoffs. Without those, an agent can optimize the visible checks by doing less.

Activation is only partially observed. The system logs that a lesson was eligible and injected, records a reached-action proxy, and the corpus experiment measures its token budget and catalogue coverage. The proxy is not adherence: a tool invocation can occur for other reasons, and a lesson can be followed without the named tool. A future version can add lightweight beacons: a lesson can name a tool or a check, and the trace can record whether the corresponding action occurred after injection. A stronger design would compare adherence rates across injected and control turns and include the result in the lesson's evidence. This is a direct application of the survey's available/reached/carried distinction.

\subsection{Selection and attribution}
The online trial/control comparison is observational. The corpus leave-one-project-out experiment adds a project-level holdout for lesson coverage, but it is a catalogue test; it does not measure a model's behavior. The user's task mix, fatigue, and the model's stochasticity can differ between arms. The system should add a placebo arm with a neutral text block of similar token length, repeated runs for a subset of tasks, and a cost ledger that includes the tokens added by lessons. The formal confidence gate in the survey is attractive, but its distribution-free bound is conservative: with small samples, only very large effects can pass. A practical study should size the benchmark for the effect it hopes to detect, or use a variance-aware paired analysis with clearly stated assumptions. The reusable-holdout literature is relevant because the loop makes repeated adaptive decisions on the same stream \citep{dwork2015reusable,cawley2010overfitting,blum2015ladder}.

\subsection{Task classification and scope}
The classifier is keyword-based, English-centric, and coarse. On the 30 hand-labeled routing requests it reaches 56.7\% accuracy; statistics, data analysis, literature, and review response are the weak classes, and several technical prompts fall into general research. A multilingual or highly technical prompt may also fall into the general class, and a single prompt may span two classes. A future version could use a small local model for classification, with the deterministic classifier as a fallback, or allow the user to pin a class. The user-level store also assumes that lessons transfer across projects. The system records task classes but not project-level transfer; a longitudinal analysis should measure whether a lesson learned in one project helps or hurts in another. An explicit project scope could be added without changing the lifecycle.

\subsection{Library governance}
The lesson library is small, but it will drift. Lessons can become stale, contradict a new project convention, or accumulate around a single noisy signal. Cross-project pooling helps a sparse family in the corpus, but it can also import a project-specific convention into an unrelated paper. The current system has caps, a time-to-live, a stale count, contradiction checks, task-class scoping, and explicit forget, but no active-set audit, no usage-based retirement beyond the success comparison, and no consolidation into skills. A future version should add an active-set report, a library-drift check, and a consolidation step that turns a stable cluster of lessons into a reviewed skill file. The skill should be versioned and admitted through the same evidence path, not written directly by the loop.

\subsection{Model and harness dependence}
The evaluation uses one plugin and one host. Lessons may be model-specific: a phrasing that helps one model may be ignored or misread by another. The system records the model in the episode but does not scope or revalidate lessons when the route changes. A future version should re-enter a lesson into trial when the model family changes, and report model--harness interactions rather than a single aggregate \citep{li2026rightfit,tian2026ceiling}. The same caution applies to host versions: a prompt or tool change can alter the effect of a lesson even when the lesson text is unchanged.

\subsection{Safety and privacy}
The lint is a baseline, not a proof. It classifies the 40-example development set perfectly, but on 20 held-out paraphrases its F1 is 0.53: six unsafe instructions passed and one safe warning was rejected. Pattern matching overfits the phrases used to write the rules. The remote judge is a veto, not a safety guarantee, and it is optional. The local store is unencrypted by default; a user who shares a machine or a backup should treat it as research data. A future version should add a local semantic filter, encryption at rest, a retention policy, a redaction audit, and an adversarial test suite that is refreshed independently of the lint rules. The symlink-safe defect benchmark should also be exercised on Windows and on network filesystems.

\subsection{What remains out of scope}
\system{} does not train weights, generate executable tools at runtime, edit the audit, or revise its own admission policy. These are deliberate exclusions, not omissions. The survey's evidence does not yet justify autonomous executable change in an academic setting, and the cost of a bad check is too high. If future work revisits executable updates, it should use the normal release path, a full regression suite, and an independent audit. The same rule applies to meta-level changes: improving the proposer or the admission policy is a code change, not a lesson.
